\chapter{Conclusion and Future Work}

\section{Conclusion}

This thesis presented a complete prototype of a zero-knowledge logistic regression (ZKLR) system implemented using GoLang and GNARK library. The core objective was to determine whether logistic regression inferences could be proven using zero-knowledge while preserving correctness, privacy, and verifiability. The results confirm that ZKLR is both feasible and practical. The prototype successfully enforces correctness by accepting all valid predictions and rejecting invalid ones, ensuring that the inference process adheres to the constraints of the underlying arithmetic circuit. Zero-knowledge properties are preserved throughout the computation, ensuring that the model weights remain private in accordance with the definitions originally formalized by Goldwasser et al. The system further demonstrates succinctness as proof size remains constant (584 bytes) and verification time is independent of the dataset size, enabled by the polynomial commitment approach used in PLONK. Finally, the prototype processes 3,000 samples with reasonable proving time, confirming the practical feasibility of real-world deployment.

\subsection{Insights Gained}

Developing this prototype led to several key insights. First, the implementation highlighted limitations in using Go for arithmetic-heavy computations. Its reliance on \texttt{big.Int} for fixed-point arithmetic introduces computational overhead in languages with native 64-bit integer support. Rust or C++ would be better suited for such workloads due to tighter control over memory, hardware-level optimizations and the ability to express fixed-point operations more efficiently. Second, although the lookup-table (LUT) method for sigmoid approximation proved effective within the zero-knowledge circuit, it contributed significantly to total constraint count. While lookup-table–based computation avoids expensive exponentiation associated with the sigmoid function~$\sigma(z) = \tfrac{1}{1+e^{-z}}$~, alternative ZKP-friendly polynomial approximations may offer better scalability. Finally, working with a single-feature, binary logistic regression model helped establish a clear roadmap for scaling the system to multivariate and multi-class architectures. This incremental design approach ensures controlled growth in circuit complexity.

\section{Future Work}

Future work aims to expand the ZKLR framework into a more general, scalable, and efficient system. One major direction is migration from Go to Rust or C++. Such a transition is expected to significantly improve performance due to native fixed-point arithmetic, reduced allocation overhead, and improved parallelization primitives. Rust’s memory model and type safety, combined with libraries such as arkworks and halo2, offer additional advantages for constructing more complex ZK circuits.

Scaling the dataset and model is another important direction. A structured multi-phase development plan would facilitate progressive expansion. In the first phase, extending to two or more features would allow validation of vectorized linear computations of the form $z = w^T x + b$~ within the ZKP environment. The second phase expanding to more features while retaining binary classification, would introduce generalized weight handling and constraint modularity. The final phase will explore multi-class logistic regression, requiring circuits capable of verifying multiple logits and a softmax-like approximation under zero-knowledge constraints. This staged evolution ensures correctness and scalability at each step.

Enhancements to circuit design forms the third major area of improvement. A more flexible circuit configuration is needed to support variable feature counts and dynamic parameterization. Improved fixed-point arithmetic routines and optimized constraint layouts would reduce proving time and memory usage. The introduction of modular subcircuits for large-scale inference will further simplify debugging and maintainability. Over time, this will move the system from a proof-of-concept to a robust, adaptable zero-knowledge inference framework.

\section{Final Remarks}

This work provides a solid foundation for zero-knowledge verification of logistic regression and demonstrates that this can be done without compromising privacy or correctness. The insights gained reveal where optimizations are necessary and provide a clear, practical roadmap for advancing toward a more scalable and flexible system. The combination of improved language support, progressive dataset expansion, and refined circuit design forms a coherent path toward a second-generation ZKLR implementation. In conclusion, this thesis establishes that zero-knowledge logistic regression is not only possible, but represents a promising direction for future research in verifiable and privacy-preserving machine learning.
